\begin{sourcebox}
\centering\itshape
Most important is vocabulary; and the vocabulary of the Russian language, which has always given intense pleasure to readers of its novels and poetry, makes a particular appeal to the systematizing mind of the mathematician.
\end{sourcebox}

\introsection{1}{Plan of the book}

The five chapters are a preparation for the Reading Selections at the end, which the student should at once examine, since learning to read such passages is in fact his whole purpose. He should note their general arrangement and content, their English titles and summaries, and any other features that catch his attention. The short selections in Parts A and B are on elementary calculus and algebra, and the longer articles in Part C are on more advanced topics.

The student has four things to learn: i) the thirty-three Cyrillic letters, Chapter I; ii) how to pronounce Russian words, Chapter II; iii) something about the endings of nouns, verbs, and especially participles, Chapter III; and most important, the mathematical vocabulary, systematically developed in Chapters IV and V. The introduction is for general orientation.

The book is intended as a ``crash course'', to be absorbed at high speed. The five chapters should be worked through rapidly, with the idea in mind that the Reading Selections will provide review. The practical exercises in Chapters I and II can be done straight ahead, but some attempt should be made to memorize the inflectional forms in Chapter III. In Chapters IV and V new words should be observed with no sustained attempt at memorization, yet sharply enough so that when they turn up again in the Readings they will not seem like total strangers; in fact, the habit of noting new material without actually memorizing it, and yet in such a way that it can later be recalled instead of being completely relearned, is basic for all rapid self-instruction. The exercises consist chiefly of short excerpts from the Readings, or from similar material, with word-for-word translations; if anyone feels that there are more exercises than he needs, let him reflect that his goal is facility; all language-learning is over-learning and the essential requirement is bulk.

Chapters I and II, on the alphabet and pronunciation, have no particular connection with mathematics and could have been shorter except for our assumption that the student will be working entirely alone. But the last three chapters, and the Readings and Glossaries, are concerned exclusively with mathematics, and almost exclusively with mathematical vocabulary.

\introsection{2}{Vocabulary by inheritance, transliteration, and loan-translation}

Almost all words in Russian and English, and in many other modern languages, fall into three classes:

\begin{description}[style=nextline,leftmargin=2.2em,labelwidth=1.7em]
\item[Class One] simple, inherited, first-level words of daily life (for the Russian words we here give only the root, in a transcription); for example, \emph{sta-} stand, \emph{lag-} lay, \emph{ber-} take, \emph{vod-} lead, \emph{nos-} carry, \emph{pis-} write, \emph{hod-} go, etc.
\item[Class Two] international words like \emph{analog}, \emph{vektor}, \emph{diskriminant}, which are transliterated from Greek or Latin, in practically the same way for all modern languages.
\item[Class Three] compound verbs (and the nouns and adjectives based on them) which are formed in conscious imitation of a Latin model (and are therefore called loan-translations) by prefixing a preposition to a verb of Class One. In this book all Latin words will be quoted in the form in which they appear in English.
\end{description}

To illustrate from English, consider the noun \emph{foresight} (native to English) and \emph{prevision} (taken over from Latin). Each of them comes from a compound verb (English \emph{foresee}, Latin \emph{previse}) which has itself been formed by prefixing a preposition (English \emph{fore}, Latin \emph{pre}) to a verb of Class One (English \emph{see}, Latin \emph{video}). Words like \emph{foresight} are rare in modern English, but in Russian they are so common that any systematic study of vocabulary must take them for its central theme. (For their connection with the ``aspect'' of Russian verbs see Chapter IV.)

In Class One some of the Russian words e.g. \emph{sta-} stand, \emph{lag-} lay, \emph{pol-} full are spelled very much like the corresponding English words, the explanation being that they are cognates; i.e. the two words, Russian and English, are actually the same word in the original Indo-European language, the common ancestor of most of the languages of India and Europe. But other pairs are spelled quite differently; e.g. \emph{ber-} take, \emph{vod-} lead, \emph{hod-} go, etc. For such Russian words, which do not come from the same Indo-European root as the English word of the same meaning, it is helpful to look for cognates in other languages, particularly Latin and Greek; e.g. \emph{ber-} (to take) is cognate with \emph{-fer} in the Latin word \emph{transfer} (to take across), and \emph{hod-} (to go) with the \emph{hod-} in the English \emph{hodometer}, formed from the Greek.

In contrast to this ``vocabulary by inheritance'' in Class One, the words in Classes Two and Three are consciously taken over from Latin or Greek, either by mere transliteration (Class Two), or else (Class Three) by loan-translation, in the following way.

Consider the Latin word \emph{circumstantia} (circumstance), which we may call a second-level word, composed of the two first-level words: \emph{circum} (around) and \emph{stantia} (stance), indicating (metaphorically) that the circumstances in which one finds oneself are the things standing around. With almost all such second-level words the original Latin represents a metaphor, which in English we have not translated but have only transliterated. But the Germans, and many other nations in Europe, including the Russians, have taken the trouble to translate the metaphor. In the German word \emph{Umstand} the two component parts \emph{um} (around) and \emph{Stand} are first-level words inherited from Indo-European. The combination is a sophisticated translation made by a German scholar, part by part, from the Latin word. The tendency to form such words was given great impetus by the printing, at Strassburg in 1466, of a German translation from the Latin of the Gutenberg Bible.

In the same way, at the time of Peter the Great (1672--1725), and as a direct result of his Westernizing influence, the Russian word \emph{ob-sto-yaniye} = \emph{circum-sta-nce} was invented, on the German model, as a ``loan-translation'' (i.e. a part-by-part imitation) of the Latin word \emph{circumstantia}. Here the prefix \emph{ob-} is a first-level preposition meaning about or around (cognate with the \emph{ob} in \emph{oblate spheroid}) and the root is the same as in the Latin (and English) words for stand. Among scientists, the most important word-builder of this kind was Lomonosov (1711--1765), who changed \emph{obstoyaniye} to its modern form \emph{obstoyatel'stvo}, with an agent-suffix \emph{-tel'-} and an abstract noun suffix \emph{-stvo}.

Similar remarks hold for countless other Russian words. In the phrase \emph{proizvodnaya proizvedeniya} derivative of a product, $d(uv)/dx$, the Latin word \emph{product}, meaning that which is brought forth (e.g., by multiplication by a factor), is formed from \emph{pro} (forth) and \emph{duc-} (lead or bring; cf. \emph{abduct}, \emph{conduct}, \emph{induct}\ldots); and the Russian imitation \emph{proizvedeniye} is formed in the same way from the preposition \emph{pro} (cognate to the Latin \emph{pro}), the preposition \emph{iz} (out of, cognate to the \emph{ex} in \emph{exit}), the verb-form \emph{ved-} meaning to lead, and the abstract noun-suffix \emph{-eniye}.

The history of the word \emph{proizvodnaya} (derivative) is very much the same. The Latin name \emph{derivata}, for the function derived from a given function by the increment process, was translated into German as \emph{Ableitung} (\emph{ab} from and \emph{leiten} lead); and then the Russians, imitating in both cases a verb meaning lead, naturally produced for the word derivative a result \emph{proizvodnaya} similar to their word \emph{proizvedeniye} of a product; (for the ``vowel gradation'' in \emph{vod-}, \emph{ved-}, see §5).

\introsection{3}{Roots and prefixes}

This process of combining a prefix (or sometimes two prefixes) with the root of a verb (e.g. \emph{pro-} forth with \emph{duc-} lead) to construct a new verb \emph{produce}, together with derived nouns and adjectives like \emph{production} and \emph{productive}, accounts for almost the entire learned vocabulary of both Russian and English.

In Russian the graphic first-level components, e.g. \emph{ved-}, \emph{vod-}, lead remain much more clearly visible than in English, so that a page of Russian prose, literary or scientific, acquires a kind of vividness that makes it very pleasant to read. Consider the sentence \emph{the points A and B coincide}. Here both the English (i.e. Latin) word \emph{co-in-cide} (\emph{co-} with, \emph{-in-} into, \emph{-cide} fall) and the Russian loan-translation \emph{so-v-pad-ayut} (\emph{so-} with, \emph{v} into, \emph{pad-} fall) involve the picturesque metaphor \emph{fall into each other}. In English the metaphor is obscured by the unfamiliar Latin form \emph{-cide} (fall), but in Russian it remains vivid, because \emph{padayut} is the lively, everyday Russian word for fall.

Since the Russian scientific vocabulary thus depends on a few prefixes and a few roots, it is natural to ask: how many, and which ones?

For the prefixes a precise answer can be given (see Chapter IV), namely nineteen, including the five already mentioned \emph{v}, \emph{iz}, \emph{ob}, \emph{pro}, \emph{so} and others like \emph{ot} out (Lat. \emph{ex}, \emph{se}) and \emph{na} on (Lat. \emph{in}). Note that in English the Latin prefix \emph{co-} is also spelled \emph{con-}, \emph{col-}, \emph{com-}, \emph{cor-}, in words like \emph{connected}, \emph{collinear}, \emph{combination}, \emph{correlate}, and that the Latin \emph{in-} (also spelled \emph{il-}, \emph{im-}, \emph{ir-}, etc.) can mean in, on, into or onto.

With respect to roots, at least 300 would be necessary for a complete reference dictionary of Russian pure mathematics, but many of them would occur very seldom and for most purposes fewer than 100 are quite sufficient; the vocabulary in Chapter V is arranged under approximately 75 roots. In order to study these roots in a systematic way we must briefly examine the history of the Indo-European language.

\introsection{4}{The Indo-European language and its descendants}

Probably about the end of the third millennium B.C. (although nothing is known here for certain) speakers of Indo-European began to spread from their original homeland, perhaps near the Black Sea, eastward into Persia and India, northward and westward into eastern and central Europe, southward into the peninsulas of Greece and Italy, and elsewhere. In every case the language was greatly affected by their migrations, becoming Sanskrit in India, Proto-Slavic in eastern Europe, Proto-Germanic in central Europe, Greek and Latin in Greece and Italy, and so forth.\corrnote{底本将 Proto-Slavic 的形成地点写作 ``in the Volga region''。这一定位并非现代历史语言学通常接受的结论；这里仅把地点改为更宽泛的 ``eastern Europe''，其余叙述保留。}

Then at various later periods these daughter-languages gave birth to granddaughters, the modern languages of India and Europe; thus English, German, Dutch, etc. are daughters of Proto-Germanic; Italian, French, Spanish, etc. are daughters of Latin; and Russian, Polish, Bulgarian etc. are daughters of Proto-Slavic. In view of the fact that Indo-European itself disappeared long before the introduction of writing, its various features can only be deduced by observing its descendants.

\introsection{5}{Vowel gradation}

The most striking of these features is the vowel gradation to be seen in sets of words like the English \emph{sing, sang, sung, song} or the Russian \emph{vod-, ved-, vad-} in \emph{proizvodnaya} derivative; \emph{proizvedeniye} product; \emph{povadka} conduct, habit. Here it is clear that all the words in a set come from the same root, which is determined by the consonants, whereas the vowels indicate parts of speech, tense, etc. The vowels are said to occur in various grades, which for several reasons it is convenient to classify as: (full) e-grade, (full) o-grade, reduced grades, and zero-grade (i.e. no vowel at all). Thus for the roots \emph{lag-} lay and \emph{ber-} take:

\begin{center}
\begin{tabular}{@{}lll@{}}
\toprule
Grade & Form & Meaning\\
\midrule
e-grade & \emph{na-leg-at'} & to lie on (e.g. a point on a line)\\
o-grade & \emph{na-log} & impost, tax\\
reduced grade & \emph{na-lag-at'} & to impose (e.g. conditions)\\
o-grade & \emph{ot-bor} & selection\\
reduced grade & \emph{ot-bir-at'} & to select\\
zero grade & \emph{br-at'} & to take\\
\bottomrule
\end{tabular}
\end{center}

\introsection{6}{Consonant variation}

Vowel gradation, i.e. variation in the vowel of a root, is inherited from the original Indo-European, and therefore occurs in both English and Russian. But certain Russian roots, i.e. those ending in $g,d,z,k,t,s,h$, also show consonant variation inherited from the mother-language Proto-Slavic, which underwent a series of historical palatalizations and related sound changes rather than one single sound-shift at about 300 A.D.\corrnote{底本把下列多种交替归为 ``the following sound-shift, perhaps about 300 A.D.''。这是过度简化；这里仅改正该历史断言，后面的交替表仍按底本保留为学习用近似。} Before front vowels (i.e. vowels equivalent to English $e$ or $i$; see the vowel-scheme in Chapter I)

\begin{enumerate}[label=\roman*)]
\item $g,d,z$ may be replaced by \emph{zh} (pronounced like the $s$ in \emph{measure});
\item $k,t,s$ may be replaced by \emph{ch}, or sometimes by \emph{shch}; and $s$ sometimes by \emph{sh};
\item $h$ may be replaced by \emph{sh}.
\end{enumerate}

Thus in modern Russian the root \emph{lag-} lay also appears in such forms as \emph{lezh-}, \emph{lozh-}; the root \emph{nos-} carry as \emph{nes-}, \emph{nosh-} etc.:

\begin{center}
\begin{tabular}{@{}ll@{}}
\emph{na-lezh-it} & it lies on (e.g. a point on a line)\\
\emph{lezh-at'} & to lie\\
\emph{v-lozh-it'} & to in-lay, im-bed\\
\emph{v-nos-it'} & to im-port, in-sert\\
\emph{v-nes-eniye} & in-ser-tion\\
\emph{so-ot-nosh-eniye} & correlation, relation, reference\\
\end{tabular}
\end{center}

The Germanic language from which English developed also underwent a consonant-shift (described by ``Grimm's Law''), consisting of the three cyclic permutations $(p,f,b)$, $(t,th,d)$, $(k,h,g)$. Thus the $t$-sound in the Indo-European word for ``three'' became ``th'' in English but remained $t$ in Russian (\emph{tri}), Latin (\emph{triangle}) and Greek (\emph{trihedral}). Similarly the Russian \emph{pol-} is the English \emph{full}, the Russian \emph{do} is the English \emph{to} etc. Moreover, it is a general feature of Indo-European languages that when such consonants are combined with $l$ or $r$ (as very often happens, in English and elsewhere; e.g., \emph{pl\ldots, br\ldots, bl\ldots}), the so-called ``liquid'' consonants $l$ and $r$ may glide from one side to the other of the vowel in the syllable. Thus Latin \emph{plenitude} and Greek \emph{plethora} correspond to Russian \emph{polnost'} and \emph{fullness}; the Latin (and borrowed English) \emph{curve} is the Russian \emph{krivaya} etc. But a systematic discussion of this metathesis of liquids, or of Grimm's Law (which has many exceptions in the incomplete form stated above) would take us too far afield. The student is merely invited, in his practical task of learning Russian vocabulary, to take note of as many examples as he personally may find interesting or helpful.

\introsection{7}{The alphabet}

Before undertaking the study of Russian words it is necessary to learn the Cyrillic letters, a task made much easier by the remarkable historical continuity connecting the Phoenician, Greek, Latin and Cyrillic alphabetic traditions.\corrnote{底本原句为 ``in the entire history of the world there has been essentially only one alphabet''。这是过强的绝对化表述；这里仅把它改为与本节实际讨论范围相符的历史连续性陈述。} Although many of the historical links are missing, the general situation can be made out rather clearly.

The first kind of writing was pictographic, when words (or groups of words) are represented by pictures. In the next stage the unit of writing is the syllable and the symbols represent not sights but sounds. This stage began with monosyllabic words, as though in English a picture of an awl were to represent first the word \emph{all} and then any syllable with a similar pronunciation. Early alphabetic writing developed during the second millennium B.C. in the eastern Mediterranean and Near East, under important Egyptian influence; Phoenician later became a major consonantal alphabetic tradition.\corrnote{底本把这一过程写成约前1700年在 Phoenicia 由 Egyptian hieroglyphs 先转为音节符号、约前1300年再转为首辅音符号的单线发展。现代文字史通常采用更复杂的西奈—黎凡特背景，因此这里对这一句作最小范围的事实性修正。} The unit of writing was now the single sound and the letters were consonants.

The advantages of this Phoenician alphabet (essentially the same family of consonantal writing as Hebrew) were so great (e.g., fewer symbols were needed) that it spread widely, including westward to Greece, where the vowels were given equal standing; e.g. the vowel $a$ (called alpha in Greek) was represented by the inherited sign \emph{aleph}, which had represented a consonant not needed in Greek. From Greece the alphabet went further westward to Southern Italy and then up to Rome, influenced in its passage by the Etruscans. Here it took on our familiar Latin form, again after some changes; e.g. $C$ was used for both /k/ and /g/ before a distinct $G$ was introduced; the Greek letter $P$ (i.e. rho = $r$) contributed to the Latin $R$, and Greek pi to Latin $P$; and the letter $Z$ was dropped from the early Latin alphabet, being later restored at the end for the transliteration of learned Greek words. Finally, the Latin alphabet, with slight further modifications, spread to France, England, America and elsewhere.\corrnote{底本把 $C$ ``replaced G as the third letter because the Etruscans did not distinguish C and G''。较准确的是早期拉丁 $C$ 曾兼表 /k/ 与 /g/，后来才创造 $G$ 加以区分；此处仅改正这一点。}

But in the meantime the Greek alphabet had also made its way directly northward from Greece. The present Russian alphabet is called \emph{Cyrillic} in honor of the Greek monk Cyril, born 827 A.D. in Salonika, a Greek city surrounded at that time by Slavic-speaking populations. The alphabet directly associated with Cyril's mission in the 860s was not Cyrillic but the Glagolitic alphabet, used in the translation of Christian texts into the language now called Old Church Slavonic. The Cyrillic script itself developed in the First Bulgarian Empire around the end of the ninth and beginning of the tenth century; its precise individual inventor is not known.\corrnote{底本称 Cyrillic ``seems to have been invented almost immediately after 863''，并猜测可能由 Clement 发明。现代研究通常把早期 Cyrillic 与第一保加利亚帝国、尤其 Preslav 文学中心联系起来，而不把个人发明者视为已定论。} Since the Christianization of Kievan Rus' in 988, the East Slavic written tradition has been deeply affected by Church Slavonic in somewhat the same way as English learned vocabulary by Church Latin.\corrnote{底本写作 ``about 950 A.D., when Christianity reached Kiev''；国家性基督教化的传统纪年为 988 年。}

In 1918, several of the original Cyrillic letters were dropped from Russian by Soviet decree, e.g. the Greek theta and iota.

With this history of the alphabet it is instructive to compare the history of chess, which is securely attested in an ancestral form in India by about the sixth century A.D. and entered Sasanian Persia in roughly the same period, later spreading northward and westward through the Islamic world into Europe.\corrnote{底本写作 ``originated in India before 300 A.D., entered Persia about 500''；现有文献证据通常把可靠记载放在约6世纪。} During the passage into Europe, the name and role of the king's counsellor or vizier changed, eventually giving rise to the European queen.\corrnote{底本进一步给出 ``vizier $\to$ virgin $\to$ lady $\to$ queen'' 的逐词讹变链；这一词源链不能视为已确立事实，故正文仅保留棋子角色演变的核心比较。} But in Russian the piece is still called \ru{ферзь} (historically from the same counsellor tradition) and is grammatically masculine. In Chapter I we shall see that most of the differences in the modern American and Russian forms of the alphabet can be traced in a similar way.
